\section{The first-crossing partition}
\label{sec:partition}

Fix $n\ge1$ and a threshold $1\le K\le n$. For $(P,Q)\in\NC_n^2$ define the
\emph{grade} after $j$ points by
\[
 k_j=\frac{h_P(j)+h_Q(j)}{2}=j-c_P(j)-c_Q(j).
\]
Since $h_P(j)\equiv h_Q(j)\equiv j\pmod2$, the grade is an integer; it starts and
ends at $0$, and each point changes it by $-1$, $0$ or $+1$ according to the two
letters $DD$, $UD/DU$, $UU$ read at that point.

\begin{definition}[LOW and HIGH sectors]
\label{def:sectors}
\lean{Meanders.FirstCrossing.Sector, Meanders.FirstCrossing.InSector}
\leanok
% Root definition: no earlier statement dependencies.
The pair $(P,Q)$ is \LOW\ if $k_j<K$ for all $0\le j\le 2n$. Otherwise let
$L$ be the least index with $k_L\ge K$ and put $u=h_P(L)$, $v=h_Q(L)$; then
$(P,Q)$ belongs to the \HIGH\ sector $(L,u,v)$.
\end{definition}

\begin{lemma}[Partition and admissible sectors]
\label{lem:partition}
\lean{Meanders.FirstCrossing.firstHeightPartition, Meanders.FirstCrossing.inSector_high_admissible, Meanders.FirstCrossing.Admissible.delta_le}
\leanok
\uses{def:sectors}
Every $(P,Q)\in\NC_n^2$ is either \LOW\ or in exactly one \HIGH\ sector. A
nonempty \HIGH\ sector $(L,u,v)$ satisfies
\[
 u+v=2K,\qquad u,v\ge1,\qquad u,v\le\min(L,2n-L),\qquad u\equiv v\equiv L\pmod2,
\]
and $\Delta:=|L-n|\le d:=n-K$. There are $O(n^{2})$ admissible triples.
\end{lemma}

\begin{proof}
\leanok
The first index $L$ with $k_L\ge K$ is unique, and since the grade moves by at
most one per point, $k_{L-1}=K-1$ and $k_L=K$, whence $u+v=2K$; the step into
$L$ must raise both heights, so the two letters at point $L-1$ are $UU$ and
$u,v\ge1$. A height after $L$ points is at most $L$, and since it must return
to $0$ in the remaining $2n-L$ points it is at most $2n-L$; parity follows from
$h(j)\equiv j$. Since $2K=u+v\le2\min(L,2n-L)$, the shorter half has length at least $K$,
which is $\Delta\le n-K$. Choosing $L$ and $u$ determines $v$.
\end{proof}

The partition is imposed on \emph{all} ordered pairs before connectivity. Each
branch will be enumerated natively, that is, by extending partial source words
letter by letter, never by listing complete matchings and filtering. The
\LOW\ branch is a height-bounded version of the classical transfer
(\cref{sec:carrier}). The point of the \HIGH\ sectors is that their metadata
$(L,u,v)$ fixes, for each of $P$ and $Q$, the number of arches crossing the cut
at $L$ and hence the exact number of closing endpoints on each side of the cut.

\begin{figure}[t]
\centering
\begin{tikzpicture}[x=0.38cm,y=0.42cm,>=Stealth]
 \draw[->] (0,0) -- (15,0) node[right] {$j$};
 \draw[->] (0,0) -- (0,4.2) node[above] {$k_j$};
 \draw[dashed] (0,3) -- (15,3) node[right] {$K$};
 \draw[thick] (0,0) -- (1,1) -- (2,1) -- (3,1) -- (4,2) -- (5,1) -- (6,1) -- (7,1) -- (8,1) -- (9,1) -- (10,1) -- (11,1) -- (12,2) -- (13,1) -- (14,0);
 \node[below] at (7,-0.3) {\small \LOW: $k_j<K$ at every cut};
\end{tikzpicture}
\hspace{0.4cm}
\begin{tikzpicture}[x=0.38cm,y=0.42cm,>=Stealth]
 \draw[->] (0,0) -- (15,0) node[right] {$j$};
 \draw[->] (0,0) -- (0,4.2) node[above] {$k_j$};
 \draw[dashed] (0,3) -- (15,3) node[right] {$K$};
 \draw[thick] (0,0) -- (1,1) -- (2,2) -- (3,3);
 \draw[thick,gray] (3,3) -- (4,2) -- (5,1) -- (6,1) -- (7,1) -- (8,0) -- (9,1) -- (10,2) -- (11,1) -- (12,0) -- (13,1) -- (14,0);
 \fill (3,3) circle (2.5pt);
 \draw[dotted] (3,0) -- (3,3);
 \node[below] at (3,-0.2) {\small $L$};
 \node[below] at (8.5,-0.9) {\small \HIGH\ sector $(L,u,v)=(3,3,3)$};
\end{tikzpicture}
\caption{Grade trajectories $k_j=(h_P(j)+h_Q(j))/2$ for two pairs of order
$n=7$ with $K=3$. Left: the \LOW\ pair
$P=UUDUDDUUDDUUDD$, $Q=UDUUDUDDUUDUDD$. Right: the \HIGH\ pair
$P=UUUDDUDDUUDDUD$, $Q=UUUDDDUDUUDDUD$, whose first hit is at $L=3$ with the
last step $UU$; only the prefix up to $L$ is subject to the first-hit test.
At odd $j$ both heights are odd, so $k_j\ge1$.}
\label{fig:partition}
\end{figure}
